\section*{Bab \ref{ch:g12:sums} --- Jumlah Peubah Acak dan Hukum Bilangan Besar}

\begin{solution}{exo:g12:sums:1}
Tiap dadu berharapan
$\frac{1 + 2 + \dots + 6}{6} = \frac{7}{2}$, sehingga menurut kelinearannya
$\E(S) = \frac72 + \frac72 = 7$. Dadunya \omterm{def:g12:sums:indep}{saling bebas}, sehingga
$\V(S) = \frac{35}{12} + \frac{35}{12} = \frac{35}{6} \approx 5.83$.
\end{solution}

\begin{solution}{exo:g12:sums:2}
$\E(X) = 50$ dan $\V(X) = 100 \times 0.25 = 25$.

Markov (karena $X \geq 0$): $\P(X \geq 75) \leq \frac{50}{75} = \frac23$.

Bienaymé--Chebyshev memberi
$\P(\abs{X - 50} \geq 25) \leq \frac{25}{25^2} = 0.04$ --- dan itu bahkan
membatasi \omterm{def:g11:prob:model}{kejadian} \emph{dua sisinya}, yang $X \geq 75$ hanya separuhnya. Chebyshev
jauh lebih tajam di sini sebab ia memanfaatkan \omterm{def:g12:randvar:exp}{ragamnya}, bukan hanya rata-ratanya.
(Nilai sebenarnya $\P(X \geq 75)$ lebih kecil daripada $10^{-6}$: jadi kedua batasnya
kasar.)
\end{solution}

\begin{solution}{exo:g12:sums:3}
$F_n$ adalah \omterm{prop:g12:sums:mean}{rata-rata contoh} $n$ peubah Bernoulli$\left(\frac12\right)$,
yang beragam $\sigma^2 = \frac14$. Batasnya
\[
\P\left(\abs{F_n - 0.5} \geq 0.05\right)
\leq \frac{1/4}{n \times 0.05^2} = \frac{100}{n}
\]
bernilai $\leq 0.05$ segera setelah $n \geq 2000$.
\end{solution}

\begin{solution}{exo:g12:sums:4}
Menurut \cref{prop:g12:sums:mean}, $\E(M_n) = 5\%$ (penganekaragaman tidak
mengubah harapan imbal hasilnya) dan
$\sigma(M_n) = \frac{20\%}{\sqrt n}$. Menuntut
$\frac{20}{\sqrt n} < 4$ memberi $\sqrt n > 5$, \emph{yaitu} $n \geq 26$.
Inilah asas \emph{penganekaragaman}: menyebarkan modal atas
risiko yang \omterm{def:g12:sums:indep}{saling bebas} \omterm{def:g12:arith:divides}{membagi} risikonya (\omterm{def:g11:stat:variance}{simpangan bakunya}) dengan $\sqrt n$
tanpa mengurangi harapan imbal hasilnya.
\end{solution}

\begin{solution}{exo:g12:sums:5}
\emph{1.} Dengan mencacah $9$ pasangan yang sama mungkin: $S$ bernilai
$2, 3, 4, 5, 6$ dengan peluang
$\frac19, \frac29, \frac39, \frac29, \frac19$. Karena itu
$\E(S) = \frac{2 + 6 + 12 + 10 + 6}{9} = 4$ dan
$\E(S^2) = \frac{4 + 18 + 48 + 50 + 36}{9} = \frac{156}{9} = \frac{52}{3}$,
sehingga $\V(S) = \frac{52}{3} - 16 = \frac43$.

\emph{2.} Untuk satu peubah: $\E(X) = 2$,
$\E(X^2) = \frac{1 + 4 + 9}{3} = \frac{14}{3}$,
$\V(X) = \frac{14}{3} - 4 = \frac23$. Maka $\E(S) = 2 + 2 = 4$ dan, menurut
\omterm{def:g12:condprob:indep}{kebebasannya}, $\V(S) = \frac23 + \frac23 = \frac43$. Nilainya sama.
\end{solution}

\begin{solution}{exo:g12:sums:6}
Dengan $Y = X$: $\V(X + Y) = \V(2X) = 4\V(X)$, sedangkan
$\V(X) + \V(Y) = 2\V(X)$. Keduanya sepakat hanya ketika $\V(X) = 0$, \emph{yaitu}
ketika $X$ bernilai tetap --- jadi keaditifannya sungguh menuntut \omterm{def:g12:condprob:indep}{kebebasan}
(di sini $X$ bergantung secara maksimal pada dirinya sendiri).
\end{solution}

\begin{solution}{exo:g12:sums:7}
Di bawah kesetimbangannya, $F_n$ adalah rata-rata $n = 1200$ peubah Bernoulli dengan
$p = \frac16$, $\sigma^2 = \frac16\cdot\frac56 = \frac{5}{36}$:
\[
\P\left(\abs{F_n - \tfrac16} \geq 0.05\right)
\leq \frac{5/36}{1200 \times 0.0025} = \frac{5}{108} \approx 0.046 .
\]
Simpangan yang teramati tepat $0.05$: yaitu \omterm{def:g11:prob:model}{kejadian} yang dihasilkan dadu setimbang
dengan peluang paling banyak $4.6\%$ --- dan karena Chebyshev sangat
berhati-hati, peluang sebenarnya jauh lebih kecil. Jadi hipotesis kesetimbangannya
tidak masuk akal; dadunya besar kemungkinan berat sebelah.
\end{solution}

\begin{solution}{exo:g12:sums:8}
\emph{1.} $p(1-p) = \frac14 - \left(p - \frac12\right)^2 \leq \frac14$.

\emph{2.} Peubah $F_n$ berharapan $p$ dan beragam
$\frac{p(1-p)}{n} \leq \frac{1}{4n}$; lalu Bienaymé--Chebyshev memberi
\[
\P\bigl(\abs{F_n - p} \geq \varepsilon\bigr)
\leq \frac{p(1-p)}{n\varepsilon^2} \leq \frac{1}{4n\varepsilon^2},
\]
yang berlaku berapa pun $p$ yang tak diketahui itu.

\emph{3.} Dengan $\varepsilon = 0.03$ dan taraf $0.05$:
\[
\frac{1}{4n(0.03)^2} \leq 0.05
\iff n \geq \frac{1}{4 \times 0.0009 \times 0.05} \approx 5556 .
\]
Jadi kira-kira $5600$ orang sudah cukup menurut batas yang kasar ini (sedangkan jawaban
klasik lewat \omterm{def:g10:numbers:approx}{hampiran} normal lebih dekat ke $1100$, lihat
\cref{ch:g12:contdist}).
\end{solution}

\begin{solution}{pb:g12:sums:1}
\textbf{1.} $V(X + Y) = V(X) + V(Y) = \frac{35}{12} +
\frac{35}{12} = \frac{35}{6}$ (menurut \omterm{def:g12:condprob:indep}{kebebasannya}).

\textbf{2.} $\E(\bar X) = 3.5$;
$\sigma(\bar X) = \frac{1.71}{\sqrt{100}} \approx 0.17$: jadi
rata-rata dari seratus dadu itu memeluk $3.5$ dalam jarak seperlima angka.

\textbf{3.} $\P(X \geq 10) \leq \frac{2}{10} = 0.2$.

\textbf{4.} $\P \leq \frac{V(\bar X)}{0.5^2} =
\frac{35/1200}{0.25} \approx 0.117$: jadi paling banyak kira-kira
$12\,\%$.

\textbf{5.} Dengan $Y = -X$: $V(X + Y) = V(0) = 0$, sedangkan
$V(X) + V(Y) = 2V(X) > 0$: jadi peubah yang berkorelasi lawan secara sempurna
saling meniadakan --- penjumlahan \omterm{def:g12:randvar:exp}{ragam} adalah hak istimewa \omterm{def:g12:condprob:indep}{kebebasan}.

\textbf{6.} $\E = \frac{18}{37} - \frac{19}{37} =
-\frac{1}{37} \approx -0.027$;
$\sigma = \sqrt{1 - \left(\frac{1}{37}\right)^2} \approx
1.00$.

\textbf{7.} $\E(G) = -2.70$, $\sigma(G) = 10 \times 1.00 =
10$: jadi hanyutannya hanya $0.27\sigma$. Derau semalam menenggelamkan
keunggulannya --- sebagian besar minoritas penjudinya (kira-kira $40\,\%$,
kata loncengnya) pulang dengan unggul, dan justru itulah yang membuat
mereka datang kembali.

\textbf{8.} Pendapatan kasinonya: harapannya $+27\,027$ euro,
\omterm{def:g11:stat:variance}{simpangan bakunya} $\approx 1\,000$: jadi labanya duduk $27$
\omterm{def:g11:stat:variance}{simpangan baku} di atas nol. Pada $27\sigma$, ``kasinonya
mungkin rugi tahun ini'' bukanlah risiko, melainkan galat pembulatan yang kecil.

\textbf{9.} $\P(\text{rugi}) = \P(G \geq 27\,027$ bagi
penjudinya$)$ $\leq \frac{10^6 \times 1}{27\,027^2} \approx
0.0014$: jadi bahkan ketaksamaan paling tumpul di buku ini menjamin
bandarnya pada $99.86\,\%$ --- dan kenyataannya jauh lebih
kuat lagi.

\textbf{10.} Tiap euro yang dipertaruhkan, kapan pun dan bagaimanapun dipilih,
berharapan imbal hasil $-\frac{1}{37}$ dari dirinya sendiri; jadi menurut kelinearannya
harapan keuntungan totalnya adalah $-\frac{1}{37} \times
(\text{total yang dipertaruhkan})$, yang negatif bagi setiap siasat yang
mempertaruhkan apa pun. Sistem yang menang akan melanggar kelinearan
harapan --- tak ada penyusunan cerdik atas taruhan buruk yang menjadikannya
taruhan baik. (Sistem pelipatduaan hanya menukar banyak kemenangan kecil dengan
kekalahan bencana yang jarang.)

\textbf{11.} Teoremanya: \emph{frekuensi} merah atas
$n$ putaran \omterm{def:g12:seq:limit}{konvergen} (dalam peluang) ke $\frac{18}{37}$. Teorema itu
tidak mengatakan apa-apa tentang putaran $n + 1$: rodanya tak punya ingatan, dan
sesudah sepuluh merah peluang merahnya tetap $\frac{18}{37}$.
Memercayai sebaliknya adalah \emph{sesat pikir penjudi}, dan
pertanyaan 12 menunjukkan apa yang sebenarnya terjadi pada deretannya.

\textbf{12.} Kelebihan $100$ merah itu tidak dibayar kembali ---
putaran berikutnya adalah salinan yang adil, dan harapan kelebihannya tetap $100$.
Tetapi $\frac{100}{10^6} = 0.0001$: yaitu seperseratus angka.
Kekeliruan sesat pikirnya dalam satu kalimat: \emph{hukum bilangan
besar mengencerkan kecelakaan masa lalu di dalam samudra percobaan baru; hukum itu
tak pernah mengirim rodanya untuk menagih utang.}

\textbf{13.} $n \geq \frac{1}{4 \times 0.05 \times 0.03^2}
\approx 5\,556$ orang.

\textbf{14.} $\frac{1.96^2 \times 0.25}{0.03^2} \approx
1\,067$: jadi lima kali lebih sedikit. Chebyshev berlaku bagi \emph{setiap}
\omterm{def:g12:randvar:rv}{distribusi} dan membayar keumumannya dengan kelonggaran; sedangkan
tetapan $1.96$ milik penjajak pendapat berasal dari bentuk lonceng yang sebenarnya
pada ayunannya, yang memusat jauh lebih keras.

\textbf{15.} Marjinnya $\propto \frac{1}{\sqrt n}$, jadi menyeparuhkannya
melipatempatkan contohnya --- kira-kira $4\,400$ orang untuk
$\pm 1.5$ angka. Ketelitian dibeli dengan harga kuadrat.

\textbf{16.} Modelnya adalah $n$ pengambilan yang \omterm{def:g12:sums:indep}{saling bebas} dengan peluang
berhasil $p$ --- dan ukuran populasinya $N$ tak pernah muncul
(sebab mengambil contoh sebagian mungil dari populasi besar yang teraduk rata
justru itulah yang disandikan \omterm{def:g12:condprob:indep}{kebebasannya}). Perumpamaan para penjajaknya: untuk
mencicipi supnya, satu sendok yang teraduk rata sudah cukup ---
entah pancinya melayani sepuluh atau sepuluh ribu orang. Kuncinya ada pada
pengadukannya --- kerangka pengambilan contoh yang bias, seperti pada jajak pendapat
di jilid sekolah menengah pertama yang menipu satu negara --- bukan pada pancinya.

\textbf{17.} $315\sqrt n < 30n \iff \sqrt n > 10.5 \iff
n > 110$: jadi di atas seratus polis hanyutannya mengalahkan
deraunya, dan tiap polis tambahan melebarkan jurangnya menurut
hukum $\sqrt n$ --- pengumpulan risiko \emph{adalah} model bisnisnya.

\textbf{18.} Bagian yang sama dan \omterm{def:g12:sums:indep}{saling bebas}: $\sigma_{\text{pf}} =
\frac{20\,\%}{\sqrt{25}} = 4\,\%$: jadi harapan imbal hasilnya sama,
tetapi ayunannya seperlima. Makan siangnya: pengurangan risiko tanpa
biaya apa pun pada harapannya --- yaitu perata-rataan keacakan yang \omterm{def:g12:sums:indep}{saling bebas}.

\textbf{19.} Jika berkorelasi sempurna: portofolionya menjadi satu saham
dalam $25$ kostum, jadi $\sigma = 20\,\%$, tanpa pengurangan sama sekali.
Penganekaragaman menyewa \emph{\omterm{def:g12:condprob:indep}{kebebasan}}; dan ketika sebuah krisis
mengorelasikan segalanya (tahun 2008: semua taruhan perumahan menjadi satu taruhan),
perlindungan $\sqrt n$ menguap justru ketika ia
diperlukan.

\textbf{20.} Jumlahnya menghanyut seperti $n$ dan berayun seperti
$\sqrt n$, sehingga rata-ratanya mengendap seperti $\frac{1}{\sqrt n}$: maka
kasinonya menabung hanyutannya atas sejuta putaran; penjajak pendapatnya
membeli $\frac{1}{\sqrt{1100}}$ derau dengan anggaran telepon;
penanggung asuransinya melampaui klaimnya sendiri pada $n > 110$; dan sang
pemodal \omterm{def:g12:arith:divides}{membagi} risikonya dengan $\sqrt{25}$. Penyalahgunanya: penjudi
yang percaya pada utang (padahal hanya ada pengenceran), dan
pemodal yang percaya pada \omterm{def:g12:condprob:indep}{kebebasan} (padahal kadang hanya ada
satu taruhan). Bentuk semesta ayunannya --- yaitu
loncengnya --- adalah bintang bab sinambungnya, dan teoremanya,
yaitu teorema limit pusat, adalah puncak peluang di jilid
universitas.

\end{solution}
